% !TEX root = document.tex
% =====================================================================
% 第2章 软件基本信息与安装 —— 所有软件类型都必写
%
% 写什么：
%   1. 软件基本信息表（三线表）：名称、版本、开发单位、完成日期、
%      运行平台、编程语言、依赖。必须和申请表 application.md 完全一致！
%   2. 系统要求：硬件（最低/推荐）+ 操作系统 + 核心依赖。
%   3. 获取与安装：安装命令、环境配置、启动方式。命令行和界面都写。
%
% 硬约束：
%   - 表中软件名、版本号必须与封面、申请表、源码页眉一致，审查会逐项核对。
%   - 界面截图上的名称必须与申请表上的软件名称一致。
%   - 每条软硬件环境描述控制在50字以内（申请表字数限制）。
%   - 表格用三线表并居中，长表格用 longtable。
% =====================================================================

\section{软件基本信息与安装配置}
\label{sec:installation}

\subsection{软件基本信息}
\label{ssec:software_info}
\begin{table}[H]
\centering
\begin{tabular}{ll}
\toprule
\textbf{项目} & \textbf{内容} \\
\midrule
软件名称 & \Name \\
软件简称 & \texttt{demo-tool}（命令行入口名） \\
版本号 & \Version{} \\
开发单位 & 示例单位（请替换为实际著作权人） \\
开发完成日期 & 2026年1月1日（请替换为实际日期） \\
运行平台 & Windows / Linux / macOS \\
编程语言 & Python 3.10+ \\
依赖库 & Streamlit \\
源代码网址 & \mbox{\url{https://example.org/your-repo}}（请替换） \\
\bottomrule
\end{tabular}
\caption{软件基本信息}
\label{tab:software_info}
\end{table}

表\ref{tab:software_info}中的软件名称、版本号与说明书封面及申请表保持一致。正式提交前，请把"示例单位""示例网址"等占位内容替换为真实信息。

\subsection{系统要求与环境准备}
\label{ssec:requirements}

\subsubsection{硬件要求}
最低配置：
\begin{itemize}
    \item CPU：2核及以上
    \item 内存：4~GB及以上
    \item 存储：100~MB可用空间
\end{itemize}
推荐配置：8核CPU、16~GB内存、SSD存储。

\subsubsection{操作系统支持}
\begin{itemize}
    \item Windows 10及以上
    \item Linux 主流发行版
    \item macOS 11及以上
\end{itemize}

\subsection{软件获取与安装}
\label{ssec:install}

\subsubsection{获取源码}
\begin{lstlisting}[language=bash,style=blockstyle]
git clone https://example.org/your-repo.git
cd your-repo
\end{lstlisting}
如无法使用Git，可在仓库页面下载源码压缩包后解压。

\subsubsection{创建运行环境}
\begin{lstlisting}[language=bash,style=blockstyle]
python -m venv .venv
source .venv/bin/activate        # Windows用 .venv\Scripts\activate
pip install -r requirements.txt
\end{lstlisting}

\subsubsection{启动方式}
命令行方式直接调用入口：
\begin{lstlisting}[language=bash,style=blockstyle]
python -m demo_tool --version
python -m demo_tool clean input.csv output.csv
\end{lstlisting}
浏览器界面方式启动服务后访问本地地址：
\begin{lstlisting}[language=bash,style=blockstyle]
streamlit run src/streamlit_app.py
\end{lstlisting}
启动后在浏览器打开\url{http://localhost:8501}即可使用界面。
